\chapter{Alignment, Safety, and Multimodality}\label{ch:alignment}

The reasoning chapters trained a specialist: a model that is graded by a verifier and rewarded
for getting a boxed answer right. Most deployed models are generalists instead. They are
expected to answer trivia, write code, chat, summarise, refuse to help with clearly harmful
requests without refusing ordinary ones, and increasingly to look at images and video. This
last chapter of Part~VI covers the parts of post-training that make a generalist out of a base
model, and how to \emph{measure} whether that worked. It then opens the model up to non-text
inputs.

For a systems engineer, three things in this chapter are worth watching. First, generalist
SFT is a data-engineering problem: mixing sources, packing a corpus that does not fit in
memory, and sharding and resuming the stream across ranks. Second, evaluation of a generalist is
a \emph{serving} workload with two phases (generate, then judge) whose cost is often dominated by
the judge model rather than the model under test. Third, multimodality changes the token budget
of a request by one to three orders of magnitude, which shows up in prefill cost, KV-cache
size, and prefix caching, where it introduces a correctness trap.

\begin{figure}[t]
\centering
\resizebox{\linewidth}{!}{%
\begin{tikzpicture}[
  st/.style={draw, rounded corners=3pt, minimum width=2.45cm, minimum height=1.0cm, align=center,
    font=\small, line width=0.6pt},
  num/.style={circle, fill=labpurple, text=white, font=\scriptsize\bfseries, inner sep=1.2pt},
  arr/.style={-{Stealth[length=2mm]}, line width=0.7pt, draw=black!65}]
\node[st, fill=fdata!18, draw=keygreen] (a) at (0,0) {Instruction\\data mixtures};
\node[st, fill=fdata!18, draw=keygreen] (b) at (3.0,0) {Packed SFT\\at scale};
\node[st, fill=fmemory!20, draw=napkin] (c) at (6.0,0) {Evaluating\\generalists};
\node[st, fill=fmemory!20, draw=napkin] (d) at (9.0,0) {The eval harness\\as a system};
\node[st, fill=warnred!12, draw=warnred] (e) at (0,-1.7) {Safety: measure\\both failure modes};
\node[st, fill=warnred!12, draw=warnred] (f) at (3.0,-1.7) {Safety training,\\guards, red teams};
\node[st, fill=fcompute!18, draw=sysblue] (g) at (6.0,-1.7) {Image encoders\\and VLMs};
\node[st, fill=fcompute!18, draw=sysblue] (h) at (9.0,-1.7) {Multimodal\\serving systems};
\foreach \n/\k in {a/1,b/2,c/3,d/4,e/5,f/6,g/7,h/8}
  \node[num] at (\n.north west) {\k};
\draw[arr] (a)--(b); \draw[arr] (b)--(c); \draw[arr] (c)--(d);
\draw[arr] (d.south) -- ++(0,-0.3) -| (e.north);
\draw[arr] (e)--(f); \draw[arr] (f)--(g); \draw[arr] (g)--(h);
\end{tikzpicture}}
\caption{Roadmap of the chapter. Stages 1--4 turn a base model into a measured generalist
assistant; stages 5--6 add the safety dimension; stages 7--8 extend the input side to images and
video and study what that does to the serving system.}
\label{fig:align:roadmap}
\end{figure}

% =====================================================================
\section{From specialist to generalist}\label{sec:align:generalist}
\covers{CS336 A5s §2; CS336 L15}

\Cref{ch:instruct} fine-tuned GPT-2 on a single Alpaca-style dataset, and
\cref{ch:rlhf,ch:dpo,ch:grpo} optimised preferences and verifiable rewards. A generalist
post-training run strings these together. The details vary between labs, but the skeleton in
\cref{fig:align:stack} is shared by most open recipes. It has a broad SFT stage on a curated
mixture that includes some safety demonstrations, then one or more preference or RL stages,
and an evaluation suite that is run after every stage and feeds back into the data mixture.

\begin{figure}[t]
\centering
\resizebox{\linewidth}{!}{%
\begin{tikzpicture}[
  box/.style={draw, rounded corners=3pt, align=center, font=\small, minimum height=1.05cm,
    line width=0.6pt},
  data/.style={box, fill=fdata!15, draw=keygreen, text width=2.9cm, font=\footnotesize},
  model/.style={box, fill=fcompute!15, draw=sysblue, text width=2.0cm},
  ev/.style={box, fill=fmemory!20, draw=napkin, text width=12.6cm, font=\footnotesize},
  arr/.style={-{Stealth[length=2mm]}, line width=0.7pt, draw=black!65}]
\node[model] (base) at (0,0) {Base model};
\node[model] (sft) at (4.2,0) {SFT model};
\node[model] (po) at (8.4,0) {Preference-tuned};
\node[model] (rl) at (12.6,0) {RL-tuned assistant};
\node[data] (mix) at (2.1,2.0) {Instruction mixture: chat, code, math, + safety demos};
\node[data] (pref) at (6.3,2.0) {Preference pairs: helpful vs harmless};
\node[data] (rlv) at (10.5,2.0) {Prompts + verifiers or reward models};
\draw[arr] (base) -- node[below,font=\scriptsize]{SFT} (sft);
\draw[arr] (sft) -- node[below,font=\scriptsize]{DPO / PPO} (po);
\draw[arr] (po) -- node[below,font=\scriptsize]{RLVR} (rl);
\draw[arr] (mix.south) -- (2.1,0.05);
\draw[arr] (pref.south) -- (6.3,0.05);
\draw[arr] (rlv.south) -- (10.5,0.05);
\node[ev] (eval) at (6.3,-2.0) {Evaluation suite after every stage: knowledge (MCQ) $\cdot$ math
(exact match) $\cdot$ chat quality (judge win rate) $\cdot$ safety (unsafe \emph{and} over-refusal)};
\foreach \m in {sft,po,rl} \draw[arr, dashed] (\m.south) -- (\m.south |- eval.north);
\draw[arr, dashed, draw=napkin] (eval.west) -| node[pos=0.75,left,font=\scriptsize,align=right]{adjust\\mixture} (-1.6,2.0) -- (mix.west);
\end{tikzpicture}}
\caption{A generalist post-training pipeline. Each stage has its own data. The evaluation suite
runs after every stage, and its results feed back into the data choices, most often the SFT
mixture. Chapter numbers: SFT in \cref{ch:instruct}, preference optimisation in
\cref{ch:rlhf,ch:dpo}, RL with verifiable rewards in \cref{ch:grpo}.}
\label{fig:align:stack}
\end{figure}

Two practical facts shape everything that follows.

\paragraph{The prompt format is part of the model.} A base model has no notion of ``user'' and
``assistant''. When you evaluate a base model zero-shot you wrap each task in a long
\emph{instruction prompt}: a system message, a description of the dialogue format, and a stop
sequence that marks where the answer ends. After SFT, the model has learned the format it was
trained with, for example an Alpaca-style \texttt{\#\#\# Instruction / \#\#\# Response} template or a
ChatML template (\cref{sec:instruct:data}). It should then be evaluated in \emph{that} format.
Comparing a base model under one wrapper with an SFT model under another is legitimate, but
you must report both wrappers, and you must never evaluate an SFT model under a template it
was not trained on and attribute the drop to the training.

\paragraph{Generalist quality is multi-dimensional.} There is no single verifier. The usual
dashboard has four axes: \emph{knowledge} (multiple-choice exams), \emph{reasoning} (short-answer
math with exact-match grading), \emph{open-ended quality} (a judge's pairwise preference against
a reference model), and \emph{safety} (behaviour on harmful and on benign-but-edgy prompts). A
change that helps one axis routinely hurts another. Heavy safety data, for example, can raise
refusals on harmless prompts and lower chat win rates.

\begin{pitfall}[Template mismatch looks like a regression]
If an SFT model trained on the template \texttt{\#\#\# Instruction:\textbackslash n\{x\}\textbackslash n\textbackslash n\#\#\# Response:\textbackslash n}
is evaluated with a ChatML prompt, it often keeps generating past the answer, invents a new
\texttt{\#\#\# Instruction} turn, or fails to stop. Parse rates collapse and the scores look as
if SFT made the model worse. Log the exact prompt string of a few examples next to every
reported number.
\end{pitfall}

% =====================================================================
\section{Instruction data for generalist models}\label{sec:align:data}
\covers{CS336 A5s §4.1}

A generalist SFT corpus is a \emph{mixture}. Typical ingredients are multi-turn chat
transcripts (often distilled from a stronger model, as in UltraChat-style corpora), task data
reformatted as instructions (summarisation, extraction, rewriting, classification), code, math
with worked solutions, and a small set of \emph{safety demonstrations}: harmful requests paired
with good refusals, plus benign requests that merely look edgy, paired with helpful answers.

\subsection{Looking at the data first}

Before writing any training code, sample a few dozen examples per source and read them. The
questions are the same for every source:
\begin{itemize}
  \item \textbf{What task is implicit?} Many ``chat'' examples are really QA, rewriting,
        brainstorming or classification. Count them, because the mixture you think you have is
        rarely the one you have.
  \item \textbf{Is the response actually good?} Distilled data inherits its teacher's verbosity,
        hedging, and errors. Look for truncated answers, refusals in non-harmful contexts, and
        canned openings repeated across thousands of examples. The model will copy these too.
  \item \textbf{Is it in-distribution for the evaluations?} If the eval set is short-answer math and the
        SFT data never shows a final numeric answer, expect parse failures.
  \item \textbf{Is it contaminated?} Benchmark questions leak into web-derived chat data.
\end{itemize}

\subsection{Mixture weights}

Given sources with $n_i$ examples, a common family of sampling distributions is
\begin{equation}
  p_i \;=\; \frac{n_i^{\alpha}}{\sum_j n_j^{\alpha}}, \qquad 0\le\alpha\le 1 .
  \label{eq:align:mix}
\end{equation}
With $\alpha=1$ you sample in proportion to size, so small sources barely register. With
$\alpha=0$ every source is equally likely, so small sources are repeated many times. Values in
between, such as $\alpha\approx0.3$--$0.7$ (the same trick multilingual pretraining uses for
low-resource languages), up-weight small but important sources such as safety data without
letting them dominate. Two numbers should be reported for every source: its sampling probability
and the \emph{number of epochs} it implies for the planned number of training examples, because
small sources that are repeated ten times or more get memorised.

\begin{lstlisting}[caption={Temperature-style mixture weights with optional per-source caps (\code{llmsys/alignment.py}).},label={lst:align:mix}]
def mixture_weights(sizes, alpha=1.0, caps=None):
    raw = {k: float(n) ** alpha for k, n in sizes.items() if n > 0}
    z = sum(raw.values())
    p = {k: v / z for k, v in raw.items()}
    # ... optional caps: clip a source's share and redistribute the excess ...
    return p

def epochs_per_source(sizes, probs, total_examples):
    return {k: probs.get(k, 0.0) * total_examples / max(1, n) for k, n in sizes.items()}
\end{lstlisting}

\begin{napkinmath}[How many times is the safety data repeated?]
Take 200{,}000 chat examples and 2{,}000 safety demonstrations, and train for 200{,}000 samples.
With $\alpha=1$, safety gets $p=2000/202000\approx1.0\%$, about 2{,}000 draws or one epoch.
With $\alpha=0.5$, $p=\sqrt{2000}/(\sqrt{200000}+\sqrt{2000})=44.7/(447.2+44.7)\approx9.1\%$,
about 18{,}200 draws, or roughly 9 epochs of the safety set. That is enough repetition to
memorise canned refusals verbatim, which shows up as an over-refusal problem
(\cref{sec:align:safety}). The fix is either a smaller $\alpha$ for only that source (a cap) or
more and more diverse safety data.
\end{napkinmath}

\begin{keyidea}[A little safety data goes a long way, too much backfires]
Public studies of safety-tuned open models (for example Bianchi et al., 2024) report that
adding a few hundred to a few thousand safety demonstrations to an otherwise ordinary SFT
mixture markedly reduces harmful compliance. Pushing the fraction much higher makes models
start refusing benign prompts that only share surface features with harmful ones. Treat the
safety fraction as a hyperparameter that is tuned against \emph{both} safety metrics of
\cref{sec:align:safety}.
\end{keyidea}

\subsection{Decontamination}

Before training, drop or flag any training example that overlaps an evaluation set. The
traditional heuristic, used for the GPT-3 era benchmarks, is a word $n$-gram overlap test with
$n\approx13$: an example is contaminated if it shares any 13-gram with any eval item. In
production the eval $n$-grams go into a Bloom filter (\cref{ch:data}) so the test streams
over billions of training documents. The book's version is deliberately simple:

\begin{lstlisting}[caption={Flagging training texts that share an $n$-gram with an eval set.},label={lst:align:decon}]
def contamination_flags(train_texts, eval_texts, n=13,
                        tokenize=lambda s: s.lower().split()):
    bank = set()
    for t in eval_texts:
        bank |= ngrams(tokenize(t), n)
    return [bool(ngrams(tokenize(t), n) & bank) for t in train_texts]
\end{lstlisting}

Overlap tests miss paraphrases and translations, so they are a lower bound on contamination.
\Cref{sec:align:evalfamilies} returns to this from the evaluation side.

% =====================================================================
\section{Packed SFT for corpora that do not fit in memory}\label{sec:align:packing}
\covers{CS336 A5s §4.2}

\Cref{sec:instruct:packing} packed examples by first-fit bin packing: load every example,
sort by length, and fill fixed-length rows, with a block-diagonal attention mask so that documents in the
same row cannot see each other. That is the right tool when the dataset fits in RAM and you
care about exact per-example semantics. A generalist corpus of hundreds of thousands to
millions of conversations, streamed from compressed JSONL shards, calls for the simpler
\emph{concat-and-chunk} layout that pretraining uses (\cref{ch:pretrain}): format each example,
append an end-of-text token, concatenate everything into one stream, and cut it into
consecutive rows of $m$ tokens.

\begin{figure}[t]
\centering
\colorlet{tkP}{fdata!35}\colorlet{tkR}{fdata!85}\colorlet{tkQ}{fcomm!30}%
\colorlet{tkS}{fcomm!80}\colorlet{tkE}{black!70}\colorlet{tkX}{black!8}%
\begin{tikzpicture}[x=0.45cm,y=0.55cm, lbl/.style={font=\scriptsize, anchor=east, align=right}]
\newcommand\tokrow[2]{\foreach \c [count=\i from 0] in {#2} {
    \fill[\c] (\i,#1) rectangle ++(1,0.85);
    \draw[white, line width=0.6pt] (\i,#1) rectangle ++(1,0.85);}}
\node[lbl] at (-0.3,5.9) {(a) pad to\\the longest};
\tokrow{6}{tkP,tkP,tkR,tkR,tkR,tkE,tkX,tkX,tkX,tkX,tkX,tkX}
\tokrow{5}{tkQ,tkQ,tkQ,tkS,tkS,tkS,tkS,tkS,tkE,tkX,tkX,tkX}
\node[lbl] at (-0.3,2.9) {(b) concat\\and chunk};
\tokrow{3}{tkP,tkP,tkR,tkR,tkR,tkE,tkQ,tkQ,tkQ,tkS,tkS,tkS}
\tokrow{2}{tkS,tkS,tkE,tkP,tkP,tkR,tkR,tkR,tkE,tkQ,tkQ,tkQ}
\node[lbl] at (-0.3,0.4) {(c) packed +\\block mask};
\tokrow{0}{tkP,tkP,tkR,tkR,tkR,tkE,tkQ,tkQ,tkQ,tkS,tkS,tkS,tkS,tkS,tkE}
\draw[decorate,decoration={brace,amplitude=3pt,mirror}] (0,-0.1) -- (6,-0.1)
  node[midway,below=3pt,font=\tiny]{doc 1 attends only to doc 1};
\draw[decorate,decoration={brace,amplitude=3pt,mirror}] (6,-0.1) -- (15,-0.1)
  node[midway,below=3pt,font=\tiny]{doc 2 attends only to doc 2};
\draw[densely dashed, warnred] (12,1.85) -- (12,4.0);
\node[font=\tiny, text=warnred, anchor=west, align=left] at (12.15,3.45) {doc 2 split\\across rows};
\begin{scope}[shift={(18.2,6.6)}]
\foreach \c/\t [count=\j from 0] in {tkP/prompt (doc 1), tkR/response (doc 1),
    tkQ/prompt (doc 2), tkS/response (doc 2), tkE/EOT delimiter, tkX/padding}{
  \fill[\c] (0,-\j*1.0) rectangle ++(0.8,0.65);
  \node[anchor=west,font=\scriptsize] at (1.0,-\j*1.0+0.33) {\t};}
\end{scope}
\end{tikzpicture}
\caption{Three layouts for SFT batches. Padding (a) wastes compute on grey tokens. Concat-and-chunk
(b) has no waste, at the cost of cross-document attention and documents split over row
boundaries. Packing with a block-diagonal mask (c, \cref{sec:instruct:packing}) has neither problem
but needs the whole dataset or a large buffer up front, plus a varlen attention kernel.}
\label{fig:align:pack}
\end{figure}

Concat-and-chunk has three consequences that you should decide on deliberately rather than
inherit by accident.

\begin{enumerate}
  \item \textbf{Cross-document attention.} Tokens of one conversation can attend to the end of
        the previous one in the same row. With an end-of-text delimiter the model learns that
        what precedes it is irrelevant, and at pretraining scale this is standard practice. For
        short SFT corpora some recipes prefer the block mask of \cref{fig:align:pack}(c); the
        stream can emit the \code{doc\_ids} needed to build it.
  \item \textbf{Documents split across rows.} A conversation that straddles a row boundary is
        trained with a truncated context in the second row. With $m=512$ and conversations of
        about 1{,}000 tokens, almost every example is split. Use $m$ at least as large as the
        typical example if you can afford it.
  \item \textbf{Which tokens carry loss.} The simplest choice trains on every token, prompt included,
        exactly like pretraining. Masking the prompt (\cref{sec:instruct:maskprompt}) focuses
        capacity on responses and is the more common choice for chat SFT. Whichever you pick,
        make it a logged flag, because it changes the loss scale and the learning curve.
\end{enumerate}

\begin{lstlisting}[caption={The streaming packer's core loop (\code{llmsys.alignment.PackedSFTStream}).},label={lst:align:stream}]
def _tokens(self):                       # yields (token, is_prompt, doc_id)
    head, tail = self.template.split("{response}")
    for d, ex in enumerate(self.docs):
        for t in self.encode(head.format(prompt=ex["prompt"])):
            yield t, True, d
        for t in self.encode(ex["response"] + tail):
            yield t, False, d
        yield self.eot_id, False, d      # the model must learn to stop

def __iter__(self):
    L, buf, row = self.seq_len, [], 0
    for item in self._tokens():
        buf.append(item)
        if len(buf) == L + 1:            # L inputs + L shifted targets
            if row >= self.skip_rows and row % self.world == self.rank:
                self.rows_emitted += 1
                yield self._make_row(buf)
            row += 1
            buf = buf[-1:]               # last token starts the next row
\end{lstlisting}

\begin{pitfall}[Tokenise the prompt and the response separately]
It is tempting to tokenise the full formatted string once and then tokenise the prompt alone to
find where the response starts. With byte-level BPE the two can disagree at the boundary: a
merge that spans the last prompt character and the first response character exists in the
full string but not in the prompt-only string, so the mask is off by one token. Encoding the
two parts separately and concatenating the ids makes the boundary exact, at the price of a
tokenisation that differs very slightly from what inference produces on the same string.
\Cref{sec:instruct:maskprompt} discusses the same trade-off for in-memory datasets.
\end{pitfall}

\subsection{Sharding and resuming the stream}

In data-parallel training every rank must see different rows, and after a crash the run must
resume \emph{exactly} where it stopped (\cref{ch:optim}). The packer handles both with two
integers. Rows are numbered globally as they are cut from the stream; rank $r$ of $W$ keeps
rows with $\text{row}\bmod W=r$; and a checkpoint stores \code{skip\_rows}, the number of global
rows already consumed. Because every rank must still tokenise the rows it discards, the cost of
this scheme is $W\times$ redundant tokenisation. That is negligible next to the forward and
backward passes for a 1B+ model, but not for a tiny one. Large pipelines pre-tokenise to binary shards
(\cref{ch:data}) and assign whole shards to ranks instead.

\begin{systemslens}[Throughput and memory budget of a generalist SFT run]
Suppose 200{,}000 conversations averaging 1{,}200 tokens, packed into $m=512$ rows with
32 rows per optimiser step: $2.4\times10^8$ tokens, $16{,}384$ tokens per step, so about
$14{,}650$ steps for one epoch. For an 8B model the training compute is
$6ND = 6\times 8.03\times10^9\times2.4\times10^8 \approx 1.16\times10^{19}$ FLOPs
(\cref{ch:pretrain}). On 8 H100s at 40\% MFU that is
$1.16\times10^{19}/(8\times0.4\times989\times10^{12})\approx 3{,}650$\,s, about an hour. The binding
constraint is memory, not compute: full fine-tuning with AdamW needs about 16 bytes per
parameter, or 128\,GB, before activations, so the model must be sharded with FSDP
(\cref{ch:sharding}). That gives 16\,GB per GPU on 8 GPUs, plus activations with checkpointing
(\cref{ch:perf}). The alternative is LoRA (\cref{sec:instruct:lora}).
\end{systemslens}

Typical full-fine-tuning SFT settings in open recipes are small learning rates (of order
$10^{-5}$) with a short linear warmup and cosine decay, weight decay around $0.1$, gradient
clipping at 1.0, and one to three epochs. Treat these as a starting point for a short learning-rate
sweep rather than as constants. \Cref{ch:optim} explains why each of them matters.

% =====================================================================
\section{Evaluating a generalist model}\label{sec:align:evalfamilies}
\covers{CS336 L12; CS336 A5s §3, §5}

\Cref{ch:eval} built the machinery for evaluating reasoning models: answer extraction,
symbolic verification, multiple-choice scoring, Elo and Bradley--Terry ratings, and rubric
judges. This section is about \emph{which} evaluations a generalist needs and how to read
them. The core difficulty is that evaluation turns an abstract construct (``is this a good
assistant?'') into a concrete metric, and every such translation leaves something out.

\subsection{What question is the evaluation answering?}

The same model can be evaluated for very different purposes: a buyer choosing between two
models for a customer-service product, a researcher measuring raw capability, a policy team
measuring risk, or a developer looking for the next thing to fix. These call for different
benchmarks. It also matters \emph{what} is being compared. Classic machine-learning evaluation
compared \emph{methods} on fixed train/test splits. Most modern leaderboards compare
\emph{systems} (model plus prompt plus tools) trained on unknown data. Speedrun-style
competitions, which fix the data and the target loss and compare wall-clock time, are the
exception that keeps evaluating methods. Whichever you do, write down the rules.

\subsection{The benchmark landscape}

\Cref{tab:align:bench} groups the evaluations a generalist is commonly judged on. The table
describes formats and grading, not scores, which go stale within months.

\begin{table}[t]
\centering\small
\begin{tabularx}{\linewidth}{@{}l>{\raggedright\arraybackslash}p{3.1cm}X@{}}
\toprule
\textbf{Family} & \textbf{Examples} & \textbf{Format, grading, and caveats}\\
\midrule
Perplexity & held-out text; LAMBADA, HellaSwag (cloze-like) & Log-likelihood on text; smooth and cheap, so scaling laws use it (\cref{ch:scalinglaws}). Conditional perplexity $p(\text{response}\mid\text{prompt})^{-1/|\text{response}|}$ ignores irrelevant tokens. Only meaningful if the model's probabilities are normalised. \\
Knowledge exams & MMLU, MMLU-Pro, GPQA, Humanity's Last Exam & Multiple choice or short answer; easy to grade; saturate as models improve, which drives harder successors. Tests knowledge more than ``understanding''.\\
Math / reasoning & GSM8K, MATH-500, ARC-AGI & Exact match after answer extraction (\cref{ch:eval}); ARC-style tasks try to separate reasoning from memorised knowledge.\\
Open-ended chat & Chatbot Arena, AlpacaEval, WildBench & Pairwise preference by humans (Arena) or an LLM judge against a reference; Elo/BT ratings (\cref{sec:eval:leaderboards}); biased towards length and style.\\
Agents & SWE-Bench, Terminal-Bench, MLE-Bench, CyBench & Multi-step tool use graded by unit tests or task success; measures model \emph{plus} scaffold; test suites can be too weak.\\
Safety & SimpleSafetyTests, HarmBench, AIR-Bench & Harmful or policy-violating requests; graded by a judge or classifier; must be paired with an over-refusal measure.\\
\bottomrule
\end{tabularx}
\caption{Families of generalist evaluations. The examples are well-known public benchmarks;
use whichever versions and splits your project specifies.}
\label{tab:align:bench}
\end{table}

\subsection{Prompting and parsing as part of the metric}

A base model asked an MMLU question in free text may answer with a letter, with the text of an option,
with a paraphrase, or with a rambling explanation. A score is only meaningful once you fix
three things:
\begin{enumerate}
  \item \textbf{A format instruction} inside the prompt, for example ``answer with a sentence of
        the form \emph{The answer is X}'' for multiple choice, or ``put the final answer in
        \texttt{\textbackslash boxed\{\}}'' for math.
  \item \textbf{A stop condition}, either an end-of-turn token or a stop string such as the next
        turn's header. Without one a base model keeps writing new questions and answering them,
        and the parser may read an answer to the wrong question.
  \item \textbf{A parser with an explicit failure value.} The fraction of unparseable outputs is a
        metric in its own right and must be reported next to accuracy. A model that ``got worse''
        may simply have stopped following the answer format.
\end{enumerate}
The book's parsers (\code{llmsys.eval.mcq}, \code{llmsys.eval.math\_verifier}) return
\code{None} on failure for exactly this reason. Use greedy decoding (temperature 0) for these
evaluations so that differences between checkpoints are not sampling noise.

\subsection{Open-ended quality: win rates and their biases}

For open-ended instructions there is no reference answer to match. The AlpacaEval recipe has a
fixed set of instructions, a fixed \emph{reference} model's outputs, and an LLM judge that
says which of two responses is better. The model's \emph{win rate} is the fraction of
instructions on which its response is preferred. Three biases need engineering
countermeasures.

\paragraph{Position bias.} Judges tend to prefer whichever response is shown first (or second,
depending on the judge). Ask twice with the order swapped and average:
\begin{lstlisting}[caption={Cancelling first-order position bias by swapping presentation order.},label={lst:align:debias}]
def debiased_preference(judge_fn, instruction, output_model, output_ref):
    p1 = judge_fn(instruction, output_model, output_ref)        # model shown first
    p2 = 1.0 - judge_fn(instruction, output_ref, output_model)  # model shown second
    return 0.5 * (p1 + p2)
\end{lstlisting}
If the judge adds a constant bonus $b$ to whichever answer comes first, the two calls return
$p+b$ and $1-(1-p+b)=p-b$, and the average recovers $p$ exactly. This doubles judge cost, which
\cref{sec:align:harness} shows is usually the dominant cost.

\paragraph{Length bias.} LLM judges, and humans, prefer longer answers even when the extra length
adds nothing. A model can climb a raw win-rate leaderboard by being verbose. AlpacaEval~2.0
introduced a \emph{length-controlled} win rate (Dubois et al., 2024): fit a regression that
explains each preference by a model term plus a term in the length difference, then report the
win rate predicted at zero length difference. A simplified version with one length covariate is
\begin{equation}
  P(\text{model preferred}) = \sigma\!\big(\theta_0 + \theta_1 \tanh(\Delta\ell/s)\big),
  \qquad \text{LC win rate} = \sigma(\theta_0),
  \label{eq:align:lc}
\end{equation}
where $\Delta\ell$ is the length difference to the reference answer and $s$ a scale (the standard
deviation of $\Delta\ell$). The $\tanh$ bounds the influence of extreme length differences.

\begin{figure}[t]
\centering
\begin{tikzpicture}
\begin{axis}[width=0.78\linewidth, height=5.4cm, xlabel={length difference to the reference, $\Delta\ell/s$},
  ylabel={$P(\text{model preferred})$}, xmin=-3, xmax=3, ymin=0, ymax=1, grid=major,
  grid style={black!10}, legend pos=south east, legend style={font=\scriptsize},
  tick label style={font=\scriptsize}, label style={font=\small}]
\addplot[sysblue, very thick, domain=-3:3, samples=100] {1/(1+exp(-(0.2+1.6*tanh(x))))};
\addlegendentry{fitted: $\theta_0=0.2$, $\theta_1=1.6$}
\addplot[napkin, very thick, dashed, domain=-3:3, samples=2] {1/(1+exp(-0.2))};
\addlegendentry{length-controlled: $\sigma(\theta_0)=0.55$}
\addplot[only marks, mark=*, mark size=1.4pt, warnred] coordinates {(1.2,0.86)};
\node[font=\scriptsize, anchor=east, text=warnred] at (axis cs:1.1,0.93) {verbose model: raw win rate $\approx0.86$};
\end{axis}
\end{tikzpicture}
\caption{The length-controlled win rate of \cref{eq:align:lc}, with illustrative parameters. A
model whose answers average $1.2s$ longer than the reference is preferred about 86\% of the
time, but the win rate predicted at equal length is only $\sigma(\theta_0)\approx0.55$.}
\label{fig:align:lc}
\end{figure}

\paragraph{Judge preferences.} A judge may prefer its own model family's style, and it can be
fooled by confident but wrong answers. Spot-check about ten examples where the model lost and ten
where it won, and record whether you agree with the judge. Rubrics and checklists
(\cref{sec:eval:judge}) make judging more reliable whether the judge is a human or a model.

\subsection{Validity: contamination, freshness, and ecological validity}

Benchmarks posted on the web end up in pretraining data. There are four practical responses:
\emph{detect} overlap (the $n$-gram test of \cref{lst:align:decon}, or statistical tests that
exploit how a memorised dataset's canonical example order is more likely than shuffled
orders); \emph{report} overlap and confidence intervals; use \emph{fresh} evaluations built
from material published after the training cutoff; and use \emph{private} evaluations (internal
code, your own documents). Separately, ask how \emph{ecologically valid} an evaluation is: exam
questions are far from real usage, and real usage logs raise privacy issues, which is why
privacy-preserving analysis of usage (for example Anthropic's Clio, Tamkin et al., 2024)
exists. Finally, benchmarks have bugs. Cleaned ``verified'' subsets (as for SWE-Bench) and
re-annotated ``platinum'' versions exist because a few percent of items in popular benchmarks
are wrong or ambiguous.

% =====================================================================
\section{The evaluation harness as a system}\label{sec:align:harness}

Evaluating a generalist means running thousands of prompts through the model under test,
running many of the outputs through a \emph{second}, often much larger, judge model, and doing this
after every checkpoint you care about. That is a batch-inference workload with two models,
and it deserves the same engineering as a serving system.

\begin{figure}[t]
\centering
\resizebox{\linewidth}{!}{%
\begin{tikzpicture}[
  box/.style={draw, rounded corners=3pt, align=center, font=\small, minimum height=1.0cm, line width=0.6pt},
  arr/.style={-{Stealth[length=2mm]}, line width=0.7pt, draw=black!65}]
\node[box, fill=fdata!15, draw=keygreen, text width=2.6cm] (suites) at (0,0) {Suites\\MCQ $\cdot$ math $\cdot$ chat $\cdot$ safety};
\node[box, fill=black!5, text width=2.4cm] (fmt) at (4.0,0) {format + key\\$h$(model, prompt,\\gen.\ config)};
\node[box, fill=fmemory!25, draw=napkin, text width=1.7cm] (cache) at (4.0,-2.1) {result\\cache (JSONL)};
\node[box, fill=fcompute!18, draw=sysblue, text width=3.0cm] (gen) at (8.4,0) {Phase 1: generate\\policy on 1 GPU\\continuous batching,\\prefix cache};
\node[box, fill=fcomm!15, draw=labpurple, text width=3.0cm] (judge) at (12.8,0) {Phase 2: judge\\70B judge, TP=2\\both orders};
\node[box, fill=black!5, text width=2.2cm] (score) at (16.6,0) {parse, score,\\CIs, report};
\draw[arr] (suites) -- (fmt); \draw[arr] (fmt) -- node[above,font=\scriptsize]{misses} (gen);
\draw[arr] (fmt.south) -- node[left,font=\scriptsize]{hits} (cache.north);
\draw[arr] (gen.south west) ++(0.4,0) |- (cache.east);
\draw[arr] (gen) -- node[above,font=\scriptsize]{outputs} (judge);
\draw[arr] (judge) -- (score);
\draw[arr] (gen.north) to[out=30,in=150] node[above,font=\scriptsize]{MCQ/math: no judge} (score.north);
\node[font=\scriptsize, text=sysblue, align=center] at (8.4,-2.1) {shared system prompt\\$\Rightarrow$ prefix-cache hits};
\node[font=\scriptsize, text=labpurple, align=center] at (12.8,-2.1) {swap models between phases\\or run on separate GPUs};
\end{tikzpicture}}
\caption{An evaluation harness as a two-phase batch-inference system. Results are cached by
(model, prompt, generation config), so re-running after a crash or re-scoring with a new parser
is free. MCQ and math suites skip the judge; chat and safety suites need it.}
\label{fig:align:harness}
\end{figure}

\begin{lstlisting}[caption={Running several suites through one batched generator with a result cache.},label={lst:align:runner}]
runner = MultiSuiteRunner(generate_fn, model_id="sft-step-14650",
                          gen_cfg={"temperature": 0.0, "top_p": 1.0},
                          cache_path="eval_cache.jsonl", batch_size=64)
mcq = Suite("knowledge", rows, format_fn=lambda ex: wrap(mcq_prompt(ex)),
            score_fn=score_letter, stops=("# Query:",), max_new_tokens=32)
report = runner.run([mcq, math_suite, chat_suite, safety_suite])
for name, r in report.items():
    print(name, r["mean"], r["ci95"], r["parse_failures"], r["examples_per_s"])
\end{lstlisting}

\Cref{lst:align:runner} shows the shape of \code{llmsys.alignment.MultiSuiteRunner}. Its
design choices are the ones from \cref{sec:eval:harnessdesign}, applied across suites:
deterministic keys that include the decoding configuration, an append-only cache that survives
crashes, stop strings applied in one place, and per-suite throughput and parse-failure counts.
The \code{generate\_fn} can be the serving engine from \cref{ch:serving}, vLLM, or a plain loop.

\begin{napkinmath}[Where does the time go in a generalist eval?]
Take an 8B model on one H100 (989\,TFLOP/s dense BF16, 3.35\,TB/s).
\begin{itemize}
  \item \emph{Knowledge MCQ}, about 14{,}000 questions of about 600 prompt tokens with a long shared system prompt,
        and a handful of output tokens. This is almost pure prefill:
        $2\times8.03\times10^9\times600\times14{,}042\approx1.35\times10^{17}$ FLOPs, or about 274\,s at 50\% of peak.
        If 150 of those tokens are a shared system prompt served from the prefix cache, prefill
        drops by a quarter, to about 205\,s.
  \item \emph{Open-ended chat}, about 800 instructions with up to 512 output tokens. This is decode-bound: each
        decode step reads about 16\,GB of weights, so it takes at least $16\times10^9/3.35\times10^{12}\approx4.8$\,ms.
        With a batch of 256, four waves of 512 steps take about 10\,s as a lower bound.
  \item \emph{The judge}: a 70B judge needs two 80\,GB GPUs for its BF16 weights. Each comparison reads about
        1{,}000 tokens (instruction plus two answers) and emits a few, and with order swapping there are $2\times805=1{,}610$
        comparisons: $2\times70.6\times10^9\times1000\times1610\approx2.3\times10^{17}$ FLOPs,
        about 230\,s on the two GPUs at 50\%.
\end{itemize}
The judge costs more than generating the answers it judges. Budget for it, cache its verdicts
by (judge, prompt) as carefully as the generations, and use batched prefill-only scoring
(\cref{sec:eval:mcq}) instead of free-text judging where the judge only needs to emit one token.
\end{napkinmath}

\begin{systemslens}[Placement: one GPU pool, two models]
On a single node there are two reasonable layouts. \emph{Phase-serial}: load the policy, generate
everything, free it, load the judge with tensor parallelism, and judge everything. This is
simple and leaves the whole node to each model in turn, but loading the judge takes minutes
(reading about 140\,GB from disk at a few GB/s). \emph{Phase-parallel}: dedicate GPUs to each model
and stream outputs from generator to judge through a queue. This overlaps the phases and
amortises loading across checkpoints, the same producer--consumer pattern as the RL systems
of \cref{ch:rlsys}. For periodic evaluation during training, keep a long-lived judge
server and send it requests over HTTP; then the trainer never pays the judge's load time.
\end{systemslens}

\begin{pitfall}[Batch size changes greedy outputs]
Greedy decoding is deterministic for a fixed batch layout, but batched kernels can reduce in a
different order for different batch sizes and padding, which flips near-tied tokens
(\cref{sec:eval:harnessdesign}). Two runs of the same checkpoint with different
\code{batch\_size} can differ by a few tenths of a point. Put the batch size in the cache key
if you need bit-exact reproducibility, and in any case report intervals rather than single
numbers: \code{MultiSuiteRunner} returns a 95\% Wilson interval per suite.
\end{pitfall}

Report throughput as examples per second per suite. It is the number that tells you whether
the next evaluation run will take ten minutes or a night, and it moves after SFT because tuned
models usually write longer answers and stop more reliably than base models.

% =====================================================================
\section{Safety: measuring both ways to fail}\label{sec:align:safety}
\covers{CS336 L12; CS336 A5s §3.4, §5.4}

``Safety'' covers a wide range of concerns: helping with clearly harmful activities, hallucinating
confidently, sycophancy, privacy leaks, and harms that depend on jurisdiction and culture. Many
of these are contextual, and some capabilities are \emph{dual-use}: a model that is good at
finding software vulnerabilities is useful both to defenders running penetration tests and to
attackers. This section takes the narrow, measurable slice that post-training targets most
directly: \emph{does the model decline to give meaningful help with clearly harmful requests,
while still helping with everything else?}

\subsection{Two error rates, not one}

Model behaviour on safety-relevant prompts has two failure modes, and they trade off against each
other:
\begin{itemize}
  \item \textbf{Unsafe compliance}: on a prompt that \emph{should} be declined, the model gives
        meaningful help. This is measured on a set of clearly harmful prompts
        (SimpleSafetyTests-style probes or larger taxonomies such as HarmBench or AIR-Bench) and
        graded by a judge model or classifier.
  \item \textbf{Over-refusal}: on a prompt that \emph{should} be answered but shares surface
        features with harmful ones (``how do I kill a Python process'', a history question about a
        war, a medical question about medication interactions), the model refuses.
\end{itemize}
A model that refuses everything has zero unsafe compliance and is useless. Report both numbers
with intervals; \code{safety\_report} in \code{llmsys.alignment} takes per-example judge labels
and does exactly that (\cref{lst:align:safety}).

\begin{figure}[t]
\centering
\begin{tikzpicture}
\begin{axis}[width=0.7\linewidth, height=5.6cm, xlabel={over-refusal rate on benign prompts},
  ylabel={unsafe-compliance rate}, xmin=0, xmax=1, ymin=0, ymax=1, grid=major, grid style={black!10},
  tick label style={font=\scriptsize}, label style={font=\small}]
\addplot[sysblue, very thick, domain=0.02:1, samples=60] {0.04/(x+0.04)};
\addplot[keygreen, very thick, dashed, domain=0.01:1, samples=60] {0.012/(x+0.012)};
\addplot[only marks, mark=*, mark size=2pt, warnred] coordinates {(0.03,0.57)};
\node[font=\scriptsize, anchor=west, text=warnred] at (axis cs:0.05,0.6) {base model};
\addplot[only marks, mark=square*, mark size=2pt, napkin] coordinates {(0.85,0.045)};
\node[font=\scriptsize, anchor=south, text=napkin] at (axis cs:0.85,0.13) {refuse-everything};
\node[font=\scriptsize, text=sysblue, anchor=west] at (axis cs:0.16,0.33) {more safety data / higher $\beta$ $\rightarrow$};
\node[font=\scriptsize, text=keygreen, anchor=west, align=left] (gl) at (axis cs:0.4,0.55) {better data: the frontier\\moves toward the origin};
\draw[-{Stealth[length=1.5mm]}, keygreen] (axis cs:0.42,0.48) -- (axis cs:0.12,0.1);
\end{axis}
\end{tikzpicture}
\caption{The safety trade-off (illustrative curves, not measurements). Turning up the amount of
safety data or the strength of safety preference training slides a model along a frontier from
the top-left to the bottom-right. The goal of better data and methods is to move the frontier towards the origin,
not just to pick a different point on it.}
\label{fig:align:safetyfront}
\end{figure}

\begin{lstlisting}[caption={A keyword refusal baseline and the two-sided safety report.},label={lst:align:safety}]
def refusal_score(text):
    head = text.strip().lower()[:200].replace(CURLY_APOSTROPHE, "'")
    return 1.0 if any(m in head for m in REFUSAL_MARKERS) else 0.0

def safety_report(records):      # record: kind, unsafe (judge), refused
    harm = [r for r in records if r["kind"] == "harmful"]
    ben = [r for r in records if r["kind"] == "benign"]
    k_unsafe = sum(bool(r["unsafe"]) for r in harm)
    k_over = sum(bool(r["refused"]) for r in ben)
    return {"unsafe_rate": k_unsafe / max(1, len(harm)),
            "unsafe_ci95": wilson_interval(k_unsafe, len(harm)),
            "over_refusal_rate": k_over / max(1, len(ben)),
            "over_refusal_ci95": wilson_interval(k_over, len(ben))}
\end{lstlisting}

The keyword baseline only inspects the opening of the reply, because many helpful answers end
with a caveat that contains ``I cannot''. It is fast and useful for triage, but it misses
partial compliance (``I can't help with that, but here is \ldots'') and polite deflections. For
reported numbers, use a judge model with an explicit question such as ``does this reply give
meaningful help towards the harmful goal?'' and validate the judge on a small hand-labelled
set, as with any judge.

\begin{napkinmath}[How many probes do you need?]
Safety probe sets are small: SimpleSafetyTests has 100 prompts. With 100 harmful prompts and 5
judged unsafe, the 95\% Wilson interval is roughly $[2\%, 11\%]$. A competing checkpoint at
8\% unsafe is \emph{not} distinguishable from it. To separate 5\% from 8\% with reasonable
power you need on the order of a thousand prompts per arm. This is why safety taxonomies grew to
thousands of prompts across hundreds of categories, and why per-category numbers from a 100-prompt
set should be read as anecdotes.
\end{napkinmath}

\subsection{How safety behaviour is trained}

Most of the machinery is the post-training machinery of earlier chapters, pointed at safety data:
\begin{itemize}
  \item \textbf{Safety SFT}: refusals and safe completions for harmful requests, and
        \emph{helpful} answers for benign look-alikes, mixed into the SFT corpus
        (\cref{sec:align:data}). The benign look-alikes are what keep over-refusal down.
  \item \textbf{Preference optimisation}: pairs in which the safe response is preferred over the unsafe one
        (helpfulness/harmlessness datasets in the style of Anthropic's HH), trained with DPO
        (\cref{ch:dpo}) or a harmlessness reward model in PPO (\cref{ch:rlhf}). Separate
        helpfulness and harmlessness reward models, combined with a constraint or a weighting,
        make the trade-off explicit.
  \item \textbf{AI feedback and constitutions}: a model critiques and revises responses against a
        written list of principles, and the revised responses or AI preference labels become
        training data (constitutional AI and RLAIF, \cref{ch:rlhf}).
  \item \textbf{Guard models at serving time}: a separate classifier screens the user input, the
        model output, or both (Llama Guard-style). It is cheap to update without retraining the
        main model, and it is defence in depth if the main model is jailbroken.
\end{itemize}

\begin{systemslens}[A guard model is a serving component with a latency budget]
An input guard adds a prefill of the user turn (small, since guards are typically 1--8B and the
inputs are short) before the main model can start, so its cost lands directly on
time-to-first-token. An output guard is harder: if it runs only on the finished response, the
user waits for the whole generation before seeing anything; if it runs on chunks of a
\emph{streamed} response, you need a policy for retracting already-sent text. Common designs run
the input guard concurrently with the main model's prefill (cancelling on a flag), classify the
output every $k$ tokens, and colocate the guard on the same GPUs as the main model to avoid a
network hop. Guard calls should be batched and prefix-cached like any other model call
(\cref{ch:serving}).
\end{systemslens}

\subsection{Red-teaming}\covers{CS336 A5s §5.5}

Red-teaming is adversarial evaluation: deliberately trying to elicit bad behaviour in order to
find failure modes before users do (Ganguli et al., 2022). The engineering side of a red-team
exercise is about method and record-keeping:
\begin{enumerate}
  \item \textbf{Scope with a taxonomy.} List the misuse categories you care about (the safety
        benchmarks above provide taxonomies), and for each one a definition of success, that is,
        what would count as meaningful help.
  \item \textbf{Vary the tactic, not just the topic.} Direct requests, role-play and fictional
        framing, multi-turn escalation, requests split into innocuous-looking parts, and
        changes of language or encoding all probe different weaknesses. Automated attacks also exist: attacker models
        that generate candidate prompts, and gradient-based search for adversarial suffixes on
        open-weight models (Zou et al., 2023), which were shown to transfer to other models.
  \item \textbf{Log every attempt}: prompt, model and version, decoding settings, outcome, and time
        spent. A red-team report says how hard the model was to break, not just whether
        it broke.
  \item \textbf{Close the loop.} Successful attacks become regression tests in the safety suite and,
        after review, training data (refusals for the attack, helpful answers for its benign
        neighbours).
\end{enumerate}

\begin{pitfall}[Handle red-team artefacts like sensitive data]
Successful attack transcripts are, by construction, examples of a model producing harmful
content. Store them with access controls, keep them out of public repositories and logs, and
review them in small doses. Using an automatic judge for the first pass is partly about scale
and partly about not exposing people to large volumes of disturbing text.
\end{pitfall}

% =====================================================================
\section{Multimodal models: everything becomes tokens}\label{sec:align:mm}
\covers{CS336 L17}

The long-term goal is an \emph{omni} model that takes any mix of text, images, audio and
video as input and can produce any of them as output. The transformer gives a uniform answer
to how: whatever the modality, turn it into a sequence of tokens, discrete ids or continuous
vectors, each carrying roughly a unit of meaning. We had to solve exactly this problem for text in
\cref{ch:text}. For images it is harder, because the natural units (pixels) are tiny and
highly redundant, and because \emph{understanding} an image (semantics) and \emph{generating}
one (fine detail) want different encodings. This section concentrates on the input side, the
vision-language model (VLM), which is the dominant deployed form, and ends with the early-fusion
alternative that can also generate images.

\subsection{Images as patch tokens}

A vision transformer (ViT; Dosovitskiy et al., 2021) cuts an $H\times W$ image into $p\times p$
patches, flattens each to a vector of $3p^2$ numbers, and projects it linearly to the model width.
A convolution with kernel and stride both equal to $p$ computes exactly this projection:
\begin{lstlisting}[caption={Patch embedding: a strided convolution is a per-patch linear map.},label={lst:align:patch}]
class PatchEmbed(nn.Module):
    def __init__(self, in_ch, patch, dim):
        super().__init__()
        self.proj = nn.Conv2d(in_ch, dim, kernel_size=patch, stride=patch)

    def forward(self, img):                  # (B, C, H, W)
        x = self.proj(img)                   # (B, dim, H/p, W/p)
        return x.flatten(2).transpose(1, 2)  # (B, N, dim), N = (H/p)(W/p)
\end{lstlisting}
At the common setting of $336\times336$ pixels with $14\times14$ patches, an image becomes
$24^2=576$ tokens. Add learned or 2-D rotary position embeddings and a stack of ordinary
(bidirectional) transformer blocks, and you have an image encoder.

\subsection{Pretraining image encoders on image--caption pairs}

Labelled image datasets are small; web pages with images and alt-text or captions are huge.
CLIP (Radford et al., 2021) trains an image encoder and a text encoder jointly so that each image
embeds close to its own caption and far from the other captions in the batch. For a batch of
$B$ pairs with normalised embeddings $\vect{u}_i$ (image) and $\vect{v}_j$ (text) and a learned
temperature $\tau$, the logits are $s_{ij}=\vect{u}_i^\top\vect{v}_j/\tau$ and the loss is a
symmetric cross-entropy, so that each image classifies its caption among $B$ and each caption classifies its image:
\begin{equation}
  \mathcal{L}_{\text{CLIP}} = -\frac{1}{2B}\sum_{i}\Big[\log\frac{e^{s_{ii}}}{\sum_j e^{s_{ij}}}
     + \log\frac{e^{s_{ii}}}{\sum_j e^{s_{ji}}}\Big].
\end{equation}
SigLIP (Zhai et al., 2023) replaces the two softmaxes by an independent binary decision per
cell, with $z_{ij}=+1$ on the diagonal and $-1$ elsewhere and a learned bias $b$:
\begin{equation}
  \mathcal{L}_{\text{SigLIP}} = -\frac{1}{B}\sum_{i,j}\log\sigma\big(z_{ij}(s_{ij}+b)\big).
\end{equation}

\begin{figure}[t]
\centering
\begin{tikzpicture}[x=0.5cm,y=0.5cm]
\begin{scope}
\foreach \i in {0,...,5} \foreach \j in {0,...,5} {
  \ifnum\i=\j \fill[keygreen!70] (\j,-\i) rectangle ++(1,-1);
  \else \fill[black!8] (\j,-\i) rectangle ++(1,-1); \fi
  \draw[white] (\j,-\i) rectangle ++(1,-1);}
\draw[sysblue, thick, rounded corners=1pt] (-0.1,-1.9) rectangle (6.1,-3.1);
\draw[napkin, thick, rounded corners=1pt] (3.9,0.1) rectangle (5.1,-6.1);
\node[font=\scriptsize, anchor=east] at (-0.2,-2.5) {row softmax};
\node[font=\scriptsize, anchor=south] at (4.5,0.2) {column softmax};
\node[font=\scriptsize, anchor=south] at (3,1.2) {rows: images, columns: captions};
\node[font=\small, anchor=north] at (3,-6.6) {CLIP: $2B$ softmaxes over $B$};
\node[font=\scriptsize] at (3,-8.1) {needs the full $B\times B$ matrix};
\end{scope}
\begin{scope}[shift={(15,0)}]
\foreach \i in {0,...,5} \foreach \j in {0,...,5} {
  \ifnum\i=\j \fill[keygreen!70] (\j,-\i) rectangle ++(1,-1);
  \else \fill[warnred!25] (\j,-\i) rectangle ++(1,-1); \fi
  \draw[white] (\j,-\i) rectangle ++(1,-1);}
\draw[labpurple, thick, dashed] (0,0) rectangle (2,-6);
\draw[labpurple, thick, dashed] (2,0) rectangle (4,-6);
\draw[labpurple, thick, dashed] (4,0) rectangle (6,-6);
\node[font=\small, anchor=north] at (3,-6.6) {SigLIP: $B^2$ independent sigmoids};
\node[font=\scriptsize, align=center] at (3,-8.5) {column blocks can be computed\\one device at a time};
\end{scope}
\end{tikzpicture}
\caption{Contrastive objectives over a batch similarity matrix. Green cells are matching
(image, caption) pairs. CLIP normalises each row and each column, which couples every cell to the
whole batch. SigLIP scores each cell on its own (positive on the diagonal, negative elsewhere),
so the matrix can be processed block by block.}
\label{fig:align:clip}
\end{figure}

\begin{systemslens}[Why the sigmoid loss scales better]
Contrastive learning wants very large batches (CLIP used 32{,}768), because more negatives per
step give a sharper signal. With data parallelism, CLIP's softmax needs every rank to see all
$B$ text embeddings and all $B$ image embeddings: an all-gather of both, then a $B\times B$ logit matrix
of which each rank uses a slice. SigLIP's loss is a plain sum over cells, so ranks can instead
pass text-embedding \emph{chunks} around a ring and accumulate the loss for one $B_{\text{local}}\times
B_{\text{local}}$ block at a time. Peak memory then no longer grows with the global batch, and the
communication overlaps with compute (the same ring pattern as ring attention, \cref{ch:distributed}).
\code{llmsys.vlm.siglip\_loss\_chunked} computes the identical value one chunk at a time; the
lab checks that it matches. The SigLIP authors also report that the sigmoid loss matches or beats
the softmax loss at moderate batch sizes, with diminishing returns beyond about 32k.
\end{systemslens}

\subsection{The LLaVA template: encoder, projector, language model}

Most open VLMs follow the template popularised by LLaVA (Liu et al., 2023): a pretrained
image encoder (CLIP or SigLIP), a small \emph{projector} that maps encoder outputs to the
language model's embedding width, and a pretrained decoder LM. The image's projected tokens are
spliced into the text-token sequence at placeholder positions (\cref{fig:align:vlm}).

\begin{figure}[t]
\centering
\resizebox{\linewidth}{!}{%
\begin{tikzpicture}[
  box/.style={draw, rounded corners=3pt, align=center, font=\small, minimum height=0.9cm, line width=0.6pt},
  tok/.style={draw=white, minimum width=0.45cm, minimum height=0.45cm, inner sep=0pt, font=\tiny},
  arr/.style={-{Stealth[length=2mm]}, line width=0.7pt, draw=black!65}]
% image grid
\foreach \i in {0,...,3} \foreach \j in {0,...,3}
  \fill[fill=fdata!\the\numexpr 20+12*\i+5*\j\relax] (\j*0.4,-\i*0.4) rectangle ++(0.38,-0.38);
\node[font=\scriptsize, anchor=north] at (0.8,-1.7) {image $\to$ $N$ patches};
\node[box, fill=fcompute!12, draw=sysblue, text width=2.0cm] (vit) at (3.6,-0.8) {vision encoder\\(ViT; frozen)};
\node[box, fill=fcomm!15, draw=labpurple, text width=1.7cm] (proj) at (6.6,-0.8) {projector\\(2-layer MLP)};
\draw[arr] (1.7,-0.8) -- (vit);
\draw[arr] (vit) -- node[above,font=\scriptsize]{$N\times d_v$} (proj);
% token sequence
\begin{scope}[shift={(8.6,-0.55)}]
\node[tok, fill=black!25] at (0,0) {s};
\foreach \k in {1,...,6} \node[tok, fill=labpurple!55] at (\k*0.47,0) {};
\foreach \k/\t in {7/W,8/h,9/a,10/t,11/?} \node[tok, fill=fdata!35] at (\k*0.47,0) {\t};
\node[font=\scriptsize, anchor=north] at (1.65,-0.35) {image slots};
\node[font=\scriptsize, anchor=north] at (4.2,-0.35) {text tokens};
\end{scope}
\draw[arr] (proj) -- node[above,font=\scriptsize]{splice} (8.35,-0.55);
\node[box, fill=fcompute!18, draw=sysblue, text width=8.5cm] (lm) at (11.4,-2.4) {decoder LM (pretrained); loss only on the response tokens};
\draw[arr] (11.4,-0.85) -- (lm.north);
\node[font=\scriptsize, align=left, anchor=west] at (0,-2.7) {stage 1: train projector only\\stage 2: projector + LM (encoder usually frozen)};
\end{tikzpicture}}
\caption{The LLaVA-style VLM. Patch features from the encoder are mapped to the LM width and
written into placeholder positions of the token sequence. The LM then treats them like any
other embeddings, so the rest of the stack, including the KV cache, is unchanged.}
\label{fig:align:vlm}
\end{figure}

\begin{lstlisting}[caption={Splicing image features into the token sequence and running a core Qwen3 decoder from embeddings (\code{llmsys/vlm.py}).},label={lst:align:splice}]
def splice_image_embeddings(tok_embeds, input_ids, image_feats, image_token_id):
    mask = input_ids == image_token_id               # (B, T) placeholder slots
    if int(mask.sum()) != image_feats.shape[0] * image_feats.shape[1]:
        raise ValueError("slot count does not match image features")
    out = tok_embeds.clone()
    out[mask] = image_feats.reshape(-1, image_feats.shape[-1]).to(out.dtype)
    return out

class TinyVLM(nn.Module):
    def forward(self, input_ids, images=None, image_feats=None, cache=None):
        x = self.lm.tok_emb(input_ids)
        if image_feats is None and images is not None:
            image_feats = self.proj(self.vision(images))
        if image_feats is not None:
            x = splice_image_embeddings(x, input_ids, image_feats, self.image_token_id)
        return qwen3_forward_embeds(self.lm, x, cache=cache)
\end{lstlisting}

The core \code{Qwen3Model.forward} always embeds token ids itself, so \code{qwen3\_forward\_embeds}
re-runs its block loop starting from embeddings. Training follows two stages:
\textbf{(1) alignment}, where the encoder and LM are frozen and only the projector is trained on
image--caption pairs, so that image tokens land somewhere the LM can read; and \textbf{(2)
instruction tuning}, where projector and LM are trained on multimodal conversations (much of it
synthesised by prompting a strong text model with captions or object annotations). Later
families add a stage that unfreezes the encoder, raise the image resolution, and add long-context
stages; the pattern of moving from easier to harder data and from frozen to unfrozen weights is shared.

\begin{figure}[t]
\centering
\begin{tikzpicture}
\begin{axis}[width=0.78\linewidth, height=5.4cm, xlabel={optimizer steps (batch 32)},
  ylabel={caption accuracy (held out)}, xmin=0, xmax=410, ymin=0, ymax=1, grid=major,
  grid style={black!10}, legend style={font=\scriptsize, at={(0.5,1.03)}, anchor=south, legend columns=2},
  tick label style={font=\scriptsize}, label style={font=\small}]
\addplot[sysblue, very thick, mark=*, mark size=1.2pt] coordinates
  {(50,0.117) (100,0.355) (150,0.516) (200,0.645) (250,0.715) (300,0.832) (350,0.836) (400,0.828)};
\addlegendentry{projector + LM from step 0}
\addplot[warnred, very thick, mark=square*, mark size=1.2pt] coordinates
  {(50,0.168) (100,0.184) (150,0.191) (200,0.188) (250,0.188) (300,0.191) (350,0.188) (400,0.188)};
\addlegendentry{projector only (stage 1 only)}
\addplot[keygreen, very thick, mark=triangle*, mark size=1.5pt] coordinates
  {(50,0.168) (100,0.184) (150,0.191) (200,0.293) (250,0.504) (300,0.598) (350,0.668) (400,0.773)};
\addlegendentry{stage 1 for 150 steps, then stage 2}
\draw[dashed, black!40] (axis cs:150,0) -- (axis cs:150,1);
\end{axis}
\end{tikzpicture}
\caption{The tiny VLM of Lab~31 learning to caption synthetic shapes (four colours times four
shapes; chance of getting both words right is $1/16$). Measured on CPU in this environment,
one seed. With a \emph{randomly initialised} LM, training the projector alone plateaus, because
there is no pretrained language ability to align to. Real VLMs start from a pretrained LM,
which is exactly why stage 1 works for them.}
\label{fig:align:vlmcurve}
\end{figure}

\subsection{Resolution, video, and where the tokens go}

Fixed-resolution encoders resize and centre-crop images (CLIP crops to $336\times336$), which
destroys small text and detail; OCR-heavy tasks suffer most. Three families of fixes are in use:
\begin{itemize}
  \item \textbf{Tiling (AnyRes)}: cut a high-resolution image into encoder-sized tiles, encode each,
        and add one downscaled overview tile. A $672\times672$ image becomes $2\times2+1=5$ tiles
        or 2{,}880 tokens; a capped grid bounds the cost of very large images (\cref{lst:align:anyres}).
  \item \textbf{Native dynamic resolution}: the encoder accepts any size directly, with 2-D rotary
        positions, and adjacent $2\times2$ patch tokens are merged before the LM so the LM sees
        a quarter as many tokens (Qwen2-VL-style).
  \item \textbf{Cross-attention resamplers}: a fixed number of learned queries (for example 256)
        cross-attend to however many patch features there are, so the LM always sees the same
        number of image tokens (Flamingo's perceiver resampler, BLIP-2's Q-Former, Qwen-VL's
        single-layer adapter).
\end{itemize}

\begin{lstlisting}[caption={AnyRes token accounting with a capped grid (\code{llmsys.vlm.anyres\_tokens}).},label={lst:align:anyres}]
def anyres_tokens(height, width, tile=336, patch=14, max_tiles=6,
                  include_base=True):
    gh, gw = math.ceil(height / tile), math.ceil(width / tile)
    while gh * gw > max_tiles:          # shrink the grid along the longer side
        if gh >= gw:
            gh -= 1
        else:
            gw -= 1
    tiles = gh * gw + (1 if include_base else 0)
    return {"grid": (gh, gw), "tiles": tiles, "tokens": tiles * (tile // patch) ** 2}
\end{lstlisting}

Video multiplies everything by the number of frames. Systems sample a few frames per second,
merge pairs of consecutive frames into one patch (3-D patches), apply the same $2\times2$ spatial
merge, and cap the total. Positions now have three axes (time, height, width). Multimodal RoPE
(M-RoPE) splits the rotary frequency bands among the three axes. Later variants interleave the
axes across bands (\texttt{t h w t h w \ldots}) rather than giving each axis a contiguous block,
so that every axis gets both low and high frequencies (\cref{ch:modern} reviews RoPE frequencies).
Explicit timestamp tokens can replace temporal positions altogether.

Two further design points from recent open VLMs are worth knowing. Injecting visual features
into \emph{several} LM layers rather than only the input (DeepStack-style) gives deeper layers
direct access to fine detail. Long video samples can dominate the token-level loss, so some
recipes rebalance with a per-sample normalisation, for example dividing by the square root of
the sample's token count rather than by the count itself.

\subsection{Early fusion: discrete image tokens}

The VLM template can read images but not write them. Early-fusion models such as Chameleon
(Chameleon Team, 2024) instead convert images into \emph{discrete} tokens with a learned
vector-quantised autoencoder (VQ-VAE: an encoder maps image patches to the nearest entries of a
codebook, and a decoder reconstructs the image from codebook entries). A $512\times512$ image
becomes about a thousand ids from a codebook of a few thousand entries, and one transformer is
trained autoregressively on interleaved text and image tokens, so generating an image is just
generating tokens. The costs are a loss of fine detail in the quantisation (OCR suffers) and
training instability. The low-entropy text tokens and high-entropy image tokens pull logits and norms
apart, and the reported fixes are the stability tools of \cref{ch:archdesign}: QK-norm and a
z-loss on the logits. Many current systems keep continuous encoders for understanding and attach
a diffusion decoder for generation.

% =====================================================================
\section{Systems for multimodal serving}\label{sec:align:mmsys}

\begin{table}[t]
\centering\small
\begin{tabular}{@{}lrrr@{}}
\toprule
\textbf{Input} & \textbf{LM tokens} & \textbf{prefill (7B)} & \textbf{KV (57\,KB/tok)}\\
\midrule
One $336^2$ image, ViT-L/14 & 576 & 8.1 TFLOP & 33\,MB\\
$672^2$ image, AnyRes $2\times2$+1 & 2{,}880 & 40 TFLOP & 165\,MB\\
$1080\times1920$ frame, AnyRes capped at 6+1 & 4{,}032 & 56 TFLOP & 231\,MB\\
60\,s video, 2 fps, $448^2$, merges $2\times2\times2$ & 15{,}360 & 215 TFLOP & 881\,MB\\
\bottomrule
\end{tabular}
\caption{Image and video token budgets (computed with \code{llmsys.vlm}). LM prefill uses
$2\times7\times10^9$ FLOPs per token; the KV column assumes a GQA model with 28 layers, 4 KV heads of
dimension 128 and BF16, about 57\,KB per token. One minute of video costs as much prefill as
roughly 30 pages of text.}
\label{tab:align:mmtokens}
\end{table}

\begin{napkinmath}[Is the vision encoder the bottleneck?]
For one $336\times336$ image a ViT-L/14 encoder (about 304M parameters) costs
$2\times3.04\times10^8\times576\approx0.35$\,TFLOP, while the 7B LM's prefill over the same 576 tokens costs
$2\times7\times10^9\times576\approx8.1$\,TFLOP. The encoder is only about 4\% of the compute, so
\emph{for images the LM prefill dominates}. Video shifts the balance: a 675M encoder over
$60\times1024=61{,}440$ patches costs about 83\,TFLOP against 215\,TFLOP of LM prefill, so the encoder is about 28\% of the total.
Token merging helps the LM but not the encoder, which still sees every patch.
\end{napkinmath}

\paragraph{Prefix caching needs image content in the key.} Prefix caching (\cref{ch:serving})
hashes chained blocks of token ids. In a VLM every image is represented by the \emph{same}
placeholder id repeated $N$ times, so two requests with identical text and different images
produce identical token ids, and a naive cache would serve one request the KV blocks computed
for the other's image. That is a silent correctness bug, not a performance one. The fix is to mix
the image's content hash into the hash of every block that contains its placeholders
(\cref{fig:align:mmhash}):

\begin{lstlisting}[caption={Block hashes that include image content (\code{llmsys.vlm.multimodal\_block\_hashes}).},label={lst:align:mmhash}]
feed, img_iter, in_run = [], iter(image_hashes), False
for tok in token_ids:
    if tok == image_token_id:
        if not in_run:                     # first slot of a new image
            current_img, in_run = next(img_iter), True
        feed.append(f"img:{current_img}")  # slot identity = image content
    else:
        in_run = False
        feed.append(str(tok))
# then chain-hash feed in blocks of block_size, exactly as for text
\end{lstlisting}

This assumes consecutive images are separated by at least one text token (real templates wrap each
image in start and end markers); otherwise two adjacent images would merge into one run.

\begin{figure}[t]
\centering
\begin{tikzpicture}[x=0.5cm,y=0.55cm,
  tk/.style={draw=white, minimum width=0.5cm, minimum height=0.5cm, inner sep=0pt, font=\tiny}]
\foreach \y/\img/\c/\lab in {0/A/fdata/request 1, -2/B/warnred/request 2} {
  \node[font=\scriptsize, anchor=east] at (-0.4,\y) {\lab};
  \foreach \k/\t in {0/s,1/{}, 2/{},3/{},4/{},5/D,6/e,7/s,8/c} {
    \ifnum\k>0 \ifnum\k<5 \node[tk, fill=\c!50] at (\k,\y) {\img}; \else \node[tk, fill=black!12] at (\k,\y) {\t}; \fi
    \else \node[tk, fill=black!12] at (\k,\y) {\t}; \fi }
  \draw[thick, sysblue] (-0.25,\y-0.35) rectangle (3.75,\y+0.35);
  \draw[thick, sysblue] (3.75,\y-0.35) rectangle (7.75,\y+0.35);
}
\node[font=\scriptsize, align=left, anchor=west] at (9.3,0.0) {token ids identical:\\\texttt{[s, img, img, img, img, D, e, \ldots]}};
\node[font=\scriptsize, align=left, anchor=west, text=warnred] at (9.3,-1.4) {ids-only hash: block 1 matches $\Rightarrow$\\request 2 reuses KV computed for image A};
\node[font=\scriptsize, align=left, anchor=west, text=keygreen] at (9.3,-2.8) {content-aware hash: \texttt{img:A} $\neq$ \texttt{img:B}\\$\Rightarrow$ miss, as required};
\end{tikzpicture}
\caption{Why multimodal prefix caching must hash image content. Both requests have the same text
and the same placeholder ids (blue boxes are KV blocks of four tokens), but different images.}
\label{fig:align:mmhash}
\end{figure}

\paragraph{Caching image embeddings.} In a multi-turn chat about an image, every turn resends
the image. With prefix caching working, the LM's KV for the image tokens is reused, but the
\emph{encoder} still runs unless its output is cached too. An LRU cache of projected features
keyed by the image's content hash, with a byte budget, removes that work (\code{ImageEmbeddingCache};
576 tokens of width 4096 in BF16 are about 4.7\,MB per image, so a 4\,GB cache holds about 850 images).

\begin{systemslens}[Separating encoding from prefill and decode]
\Cref{ch:serving} disaggregated prefill from decode because they stress different resources.
Multimodal serving adds a third stage with its own profile: image decoding and resizing on the
CPU, then a bidirectional encoder over many patches on the GPU. Some serving stacks therefore
run encoders as a separate pool (sometimes called encode--prefill--decode or EPD disaggregation)
that ships projected features to the prefill workers. This lets the encoder pool scale with the image traffic
rather than the text traffic. On the training side, the classic bottleneck is the data loader:
JPEG decoding and augmentation for thousands of images per second needs many CPU workers or
GPU decoding, and samples with very different token counts (one image versus a 60-second clip)
need the same length-aware packing and bucketing as text (\cref{sec:align:packing,ch:classify}).
\end{systemslens}

% =====================================================================
\section{Summary}\label{sec:align:summary}
\begin{itemize}
  \item A generalist post-training run is a pipeline of SFT on a curated \emph{mixture}, then
        preference or RL stages, with a multi-axis evaluation suite after every stage that feeds back
        into the data. Report sampling probability \emph{and} epochs per source; decontaminate
        against the eval sets.
  \item For corpora that do not fit in memory, stream, concatenate and chunk, with explicit
        decisions about prompt masking and cross-document attention, separate tokenisation of prompt
        and response, rank sharding, and a resumable cursor.
  \item Evaluation needs a format instruction, a stop condition, and a parser with an explicit
        failure value. Parse-failure rate and throughput are metrics. Open-ended win rates need
        position debiasing and length control; all numbers need intervals.
  \item The harness is a two-phase batch-inference system in which the judge often costs more
        than generation. Cache by (model, prompt, config), exploit shared-prefix caching, and plan
        the placement of policy and judge.
  \item Safety has two error rates, unsafe compliance and over-refusal, that trade off. Measure
        both, train with paired safety and benign look-alike data, add guard models as serving
        components, and red-team with a taxonomy, varied tactics, and logs.
  \item VLMs splice projected image features into the token stream. Images and video multiply the token
        count by orders of magnitude, LM prefill usually dominates encoder cost, and prefix
        caching must hash image content to stay correct.
\end{itemize}

% =====================================================================
\section{Exercises}\label{sec:align:exercises}

\begin{exercise}[Mixture arithmetic]
Sources have sizes chat $=500{,}000$, code $=50{,}000$, math $=20{,}000$, safety $=3{,}000$. For
$\alpha\in\{1, 0.5, 0.3\}$ compute each source's sampling probability and the number of epochs of
each source in a run of 600{,}000 samples. Which $\alpha$ keeps safety under 5 epochs while giving it
at least 1\% of samples? Then use \code{mixture\_weights(..., caps=...)} to get the same effect at
$\alpha=0.5$.
\end{exercise}

\begin{exercise}[Packing waste]
Conversations have lengths drawn uniformly from 200 to 1{,}800 tokens. For a row length
$m\in\{512, 2048\}$, estimate (a) the fraction of compute wasted by padding each conversation to the
next multiple of $m$, and (b) the fraction of conversations split across at least one row boundary
under concat-and-chunk. What does this suggest about choosing $m$ for SFT?
\end{exercise}

\begin{exercise}[Boundary tokenisation]
Using the byte-level BPE tokenizer of \cref{ch:text} trained on any text, find a (prompt, response)
pair for which \code{encode(prompt + response)} is not equal to
\code{encode(prompt) + encode(response)}. Which tokens differ? Explain why the prompt mask computed
from \code{len(encode(prompt))} would then be off by one.
\end{exercise}

\begin{exercise}[Debiasing a judge]
Suppose a judge returns $\min(1, p + b)$ for the answer shown first, where $p$ is the true preference
and $b>0$ a bias. Show that \code{debiased\_preference} recovers $p$ exactly when $p+b\le1$ and
$1-p+b\le1$. What does it return when $p+b>1$? Propose a fix.
\end{exercise}

\begin{exercise}[Length control]
Generate synthetic preferences where the true quality advantage is $\theta_0=0$ and the judge
responds to length with $\theta_1=2$. Fit \code{length\_controlled\_win\_rate} for 50, 200 and
800 examples, repeating each 20 times. How biased and how noisy is the LC estimate at each size? What happens if
the model's lengths are almost always longer than the reference (little overlap around $\Delta\ell=0$)?
\end{exercise}

\begin{exercise}[Eval budget]
Redo the napkin math of \cref{sec:align:harness} for a 1.7B policy and an 8B judge on one H100,
assuming MCQ prompts of 400 tokens. Which phase dominates now? At what judge size does judging cost the same
as generation?
\end{exercise}

\begin{exercise}[Safety statistics]
You have 100 harmful and 100 benign probes. Checkpoint A has 7 unsafe and 4 over-refusals; B has 3 unsafe
and 12 over-refusals. Compute Wilson intervals for all four rates. Can you claim B is safer? Design a
probe-set size that would let you detect a 4-point difference in unsafe rate with 80\% power.
\end{exercise}

\begin{exercise}[Guard-model latency]
An input guard of 1B parameters screens prompts of 300 tokens on the same H100 as an 8B main model. Estimate
the added time-to-first-token (a) if the guard runs before the main prefill, and (b) if it runs
concurrently. What does option (b) cost when the guard flags a request?
\end{exercise}

\begin{exercise}[CLIP versus SigLIP memory]
With a global batch of 32{,}768 over 64 data-parallel ranks and 1024-dimensional embeddings in BF16,
compute per-rank memory for (a) the gathered text and image embeddings and the local slice of the logit matrix
under CLIP, and (b) one logit block at a time under ring-chunked SigLIP. Verify numerically with
\code{siglip\_loss\_chunked} that chunking does not change the loss.
\end{exercise}

\begin{exercise}[Multimodal prefix caching]
Construct two requests that would collide under an ids-only prefix cache, and show with
\code{multimodal\_block\_hashes} that they do not collide with content-aware hashing. Then construct two
\emph{adjacent} images with no text between them and show the failure mentioned in
\cref{sec:align:mmsys}. Fix it by inserting image start/end markers.
\end{exercise}

% =====================================================================
\section{Lab 31: Generalist plumbing, evaluation at scale, and a tiny VLM}\label{sec:align:lab}

\begin{labbox}{31}{Generalist plumbing, evaluation at scale, and a tiny VLM}
\labdir{lab31_alignment}
\labsection{Goals}
Build the data, evaluation and multimodal components of this chapter against the book's
\code{llmsys} API, check them against the library versions, and train a tiny vision-language model
on CPU.

\labsection{Setup}
\begin{lstlisting}[style=shell]
cd code
LAB_IMPL=starter  python -m pytest labs/lab31_alignment      # grade your work
LAB_IMPL=solution python -m pytest labs/lab31_alignment      # reference
python labs/lab31_alignment/run.py vlm --steps 400           # CPU, minutes
\end{lstlisting}

\labsection{Tasks}
\begin{enumerate}
  \item \code{mixture\_weights}: $p_i\propto n_i^\alpha$ (\cref{eq:align:mix}).
  \item \code{pack\_stream}: streaming concat-and-chunk with separate prompt/response tokenisation,
        EOT delimiters, one-token row overlap and prompt masking. The tests compare it with
        \code{PackedSFTStream} and check that the rows re-assemble into the token stream.
  \item \code{debiased\_preference}: cancel position bias (the test uses a judge with a known bias).
  \item \code{safety\_report}: unsafe-compliance and over-refusal rates with Wilson intervals.
  \item \code{splice\_image\_embeddings}: write projected features into placeholder slots; reject a
        mismatched slot count.
  \item \code{clip\_loss} and \code{siglip\_loss}: the two contrastive objectives.
  \item \code{anyres\_tokens}: tiling with a capped grid and an overview tile.
  \item \code{multimodal\_block\_hashes}: content-aware chained block hashes.
  \item \code{ImageEmbeddingCache.get\_or\_compute}: an LRU cache of image features with a byte budget.
\end{enumerate}
The synthetic task (four colours times four shapes on $8\times8$ images, two-word captions), the model
builder and the stage-aware training loop are provided. A test checks that stage~1 leaves the LM and
encoder untouched and that stage~2 reduces the caption loss.

\labsection{Measurements}
\begin{itemize}
  \item Reproduce \cref{fig:align:vlmcurve} with \code{run.py vlm}. Then initialise the LM by first
        training it as a text-only model on the caption grammar, and rerun the stage-1-only arm: does
        aligning the projector to a \emph{pretrained} LM now work?
  \item With a real checkpoint (Qwen3-0.6B) and your own small JSONL suites, run
        \code{run.py eval} twice and confirm the second run is served entirely from the cache. Record
        examples per second for each suite and the parse-failure count.
  \item Time \code{PackedSFTStream} on a 100\,MB JSONL file with a fast tokenizer and compare tokens per second
        with your trainer's tokens per second (\cref{ch:optim}). Is the data path a bottleneck at 1 rank? At 8?
\end{itemize}

\labsection{Deliverables}
All tests passing; the VLM stage-ablation plot with your extra pretrained-LM arm; a one-page evaluation
report for one checkpoint with, per suite, the score with 95\% interval, parse failures, throughput, and
for chat and safety the judge configuration and both orders' agreement rate.

\labsection{Stretch goals}
Add 2-D rotary positions to \code{TinyViT}; implement a resampler projector (16 learned queries
cross-attending to the patch features) and compare the token count and accuracy against the MLP projector;
plug \code{multimodal\_block\_hashes} into the prefix cache of the serving engine from \cref{ch:serving}.
\end{labbox}
